\section*{Hoofdstuk \ref{ch:b2:proba} --- Kansrekening op aftelbare ruimten}

\begin{solution}{exo:b2:proba:1}
Wegens de symmetrie: de trekvolgorde brengt een uniform
willekeurige onderlinge volgorde op de ballen $1$ en $2$ voort, dus
$\P(1 \text{ vóór } 2) = \frac12$. Formeel: de posities van de
ballen $1$ en $2$ in een trekrij verwisselen is een bijectie van de
(even waarschijnlijke) uitkomsten die de \omterm{def:b2:proba:space}{gebeurtenis} met haar
complement verwisselt. Onder de ballen $1, \dots, k$: de onderlinge
volgorde van deze $k$ ballen is uniform over de $k!$
rangschikkingen, en bal $1$ staat in $(k-1)!$ daarvan vooraan:
kans $\frac{(k-1)!}{k!} = \frac1k$.
\end{solution}

\begin{solution}{exo:b2:proba:2}
$\frac{1}{k(k+1)} = \frac1k - \frac1{k+1}$, dus telescopeert
$\sum_{k\geq1} p_k$ tot $1$: een \omterm{def:b2:proba:space}{kansmaat}. Even uitkomsten:
\[
\P(2\N^*) = \sum_{j=1}^{\infty}\frac{1}{2j(2j+1)}
= \sum_{j=1}^{\infty}\Bigl(\frac{1}{2j} - \frac{1}{2j+1}\Bigr)
= \frac12 - \frac13 + \frac14 - \frac15 + \cdots
\]
Dit is de alternerende harmonische reeks zonder haar eerste term en
met omgekeerde tekens: omdat $\ln 2 = 1 - \frac12 + \frac13 -
\frac14 + \cdots$ (\cref{ch:b2:series}), is
\[
\P(2\N^*) = -\bigl(\ln 2 - 1\bigr) = 1 - \ln 2 \approx 0.307 .
\]
\end{solution}

\begin{solution}{exo:b2:proba:3}
Zij $S$ = ziek en $+$ = positieve test. Bayes
(\cref{thm:b2:proba:bayes}) met de partitie $\{S, S^c\}$:
\[
\P(S \mid +)
= \frac{0.99 \times 10^{-4}}
       {0.99 \times 10^{-4} + 0.01 \times 0.9999}
= \frac{0.000099}{0.000099 + 0.009999}
\approx 0.0098 ,
\]
onder $1\%$. Hoewel de test ``$99\%$ nauwkeurig'' is, laat een
positieve uitslag je met ongeveer $99\%$ kans gezond: de
vals-positieven onder de enorme gezonde meerderheid overspoelen de
echte positieven uit de piepkleine zieke minderheid. Screeningtests
voor zeldzame aandoeningen moeten altijd door deze berekening met
het basispercentage worden gelezen.
\end{solution}

\begin{solution}{exo:b2:proba:4}
Puntsgewijs op $\Omega$: $\omega \in \bigcup A_i$ dan en slechts
dan als een zekere factor $1 - \mathbf{1}_{A_i}(\omega)$ verdwijnt,
dus is
\[
\mathbf{1}_{\bigcup A_i}
= 1 - \prod_{i=1}^n\bigl(1 - \mathbf{1}_{A_i}\bigr)
= \sum_{\emptyset \neq J \subseteq \{1,\dots,n\}}
(-1)^{\abs J + 1}\prod_{i \in J}\mathbf{1}_{A_i} ,
\]
door het product uit te werken en de $1$ over te brengen. Nu is
$\prod_{i\in J}\mathbf{1}_{A_i} = \mathbf{1}_{\bigcap_{i \in J}
A_i}$, en sommeren tegen de gewichten $\P(\{\omega\})$ ---
geoorloofd: eindig veel begrensde termen, elke familie \omterm{def:b2:series:summable}{sommeerbaar}
--- maakt van elke indicator de kans op haar \omterm{def:b2:proba:space}{gebeurtenis}, wat de
formule geeft.
\end{solution}

\begin{solution}{exo:b2:proba:5}
Zij $A_i$ = ``brief $i$ zit in de juiste envelop''. Voor $J$ van
grootte $k$ is $\P\bigl(\bigcap_{i\in J}A_i\bigr) =
\frac{(n-k)!}{n!}$ (leg $k$ brieven vast, permuteer de rest).
Volgens de in- en uitsluiting is
\[
\P\Bigl(\bigcup A_i\Bigr)
= \sum_{k=1}^n (-1)^{k+1}\binom nk \frac{(n-k)!}{n!}
= \sum_{k=1}^n \frac{(-1)^{k+1}}{k!} ,
\]
dus
\[
\P(\text{geen overeenkomst})
= 1 - \P\Bigl(\bigcup A_i\Bigr)
= \sum_{k=0}^{n}\frac{(-1)^k}{k!}
\xrightarrow[n\to\infty]{} e^{-1} \approx 0.368 .
\]
Precies één overeenkomst: een permutatie met precies één vast punt
wordt bepaald door de keuze van de vaste brief ($n$ manieren) en
een \emph{derangement} (rangschikking zonder overeenkomst) van de
andere $n - 1$; schrijven wij $D_{n-1} =
(n-1)!\sum_{k=0}^{n-1}\frac{(-1)^k}{k!}$ voor het aantal
derangementen (het eerste deel, met $(n-1)!$ geschaald), dan is
\[
\P(\text{precies één overeenkomst})
= \frac{n\,D_{n-1}}{n!}
= \frac{D_{n-1}}{(n-1)!}
= \sum_{k=0}^{n-1}\frac{(-1)^k}{k!}
\xrightarrow[n\to\infty]{} e^{-1} :
\]
in de limiet zijn ``geen overeenkomst'' en ``precies één
overeenkomst'' even waarschijnlijk, elk met kans $e^{-1}$.
\end{solution}

\begin{solution}{exo:b2:proba:6}
Stel voorwaardelijk op het begin (kettingregel /
\cref{thm:b2:proba:bayes}):
\begin{itemize}
\item eerste worp M (kans $1 - p$): het spel begint opnieuw; langer
duren dan $n$ betekent van daar af langer duren dan $n - 1$:
bijdrage $(1-p)\,q_{n-1}$;
\item eerste worpen KM (kans $p(1-p)$): opnieuw beginnen na twee
worpen: bijdrage $p(1-p)\,q_{n-2}$;
\item eerste worpen KK: het spel is geëindigd (binnen $n$ worpen,
$n \geq 2$): bijdrage $0$.
\end{itemize}
Bijgevolg is $q_n = (1-p)q_{n-1} + p(1-p)q_{n-2}$. De
karakteristieke vergelijking $r^2 = (1-p)r + p(1-p)$ heeft
nulpunten
\[
r_\pm = \frac{(1-p) \pm \sqrt{(1-p)^2 + 4p(1-p)}}{2},
\]
met $\abs{r_\pm} < 1$: inderdaad voldoet de veelterm $\chi(r) = r^2
- (1-p)r - p(1-p)$ aan $\chi(1) = 1 - (1-p) - p(1-p) = p^2 > 0$ en
$\chi(-1) = 1 + (1-p) - p(1-p) > 0$, terwijl $\chi(0) = -p(1-p) <
0$: één nulpunt in $\intoo{-1}{0}$, één in $\intoo{0}{1}$. Dus $q_n
= \alpha r_+^n + \beta r_-^n \to 0$. De \omterm{def:b2:proba:space}{gebeurtenissen} ``het spel
duurt langer dan $n$'' dalen naar ``het spel eindigt nooit''; de
monotone \omterm{def:b2:metric:continuity}{continuïteit} (\cref{thm:b2:proba:continuity}) geeft
$\P(\text{eindigt nooit}) = \lim q_n = 0$: het spel eindigt bijna
zeker.
\end{solution}

\begin{solution}{exo:b2:proba:7}
\emph{$\P(R_n) = 1/n$:} onder de eerste $n$ trekkingen is elk van
de $n$ onderlinge posities van de laatste trekking even
waarschijnlijk (uniformiteit van de onderlinge volgorde), en $R_n$
is de \omterm{def:b2:proba:space}{gebeurtenis} dat zij de grootste is: kans $1/n$.

\emph{Oneindig veel records:} $\sum_n \P(R_n) = \sum 1/n = \infty$
en de $R_n$ zijn \omterm{def:b2:proba:independence}{onafhankelijk} (aangenomen), dus geeft
Borel--Cantelli~2 (\cref{thm:b2:proba:borelcantelli}) dat
$\P(\limsup R_n) = 1$: records houden bijna zeker nooit op --- maar
zij dunnen logaritmisch uit.

\emph{Opeenvolgende records:} wegens de \omterm{def:b2:proba:independence}{onafhankelijkheid} is
\[
\sum_n \P(R_n \cap R_{n+1})
= \sum_n \frac{1}{n(n+1)} < \infty ,
\]
dus is Borel--Cantelli~1 van toepassing: bijna zeker wordt slechts
eindig vaak een record onmiddellijk door een ander record gevolgd.
De twee helften van het lemma werken samen: oneindig veel records,
maar (bijna zeker) uiteindelijk nooit twee op rij.
\end{solution}

\begin{solution}{exo:b2:proba:8}
\emph{Eerste deel:} $\sum \P(A_n) = \sum\frac{1}{n+1} = \infty$ met
\omterm{def:b2:proba:independence}{onafhankelijkheid}: Borel--Cantelli~2 geeft $\P(\limsup A_n) = 1$.
Elke afzonderlijke $A_n$ is steeds onwaarschijnlijker, en toch
behoort bijna elke $\omega$ tot oneindig veel ervan.

\emph{Tegenvoorbeeld zonder \omterm{def:b2:proba:independence}{onafhankelijkheid}:} neem $\Omega =
\N^*$ met de gewichten $p_k = \frac{1}{k(k+1)}$ van
\cref{exo:b2:proba:2}, en $A_n = \{k \in \N^* : k \geq n\}$. Dan is
\[
\P(A_n) = \sum_{k \geq n}\Bigl(\frac1k - \frac1{k+1}\Bigr)
= \frac1n ,
\qquad
\sum_n \P(A_n) = \infty ,
\]
maar de $A_n$ zijn dalend, dus $\limsup_n A_n = \bigcap_n A_n =
\emptyset$: $\P(\limsup A_n) = 0$. De divergentie van $\sum\P(A_n)$
alleen waarborgt niets wanneer de \omterm{def:b2:proba:space}{gebeurtenissen} zich op een
krimpend deel van de ruimte opstapelen --- de \omterm{def:b2:proba:independence}{onafhankelijkheid} is
wat die samenzwering verbiedt.
\end{solution}

\begin{solution}{exo:b2:proba:9}
Met $q = 1 - p$ valt de eerste kop op rang $2j + 1$ met kans
$q^{2j}p$, dus
\[
\P(\text{oneven rang}) = \sum_{j\geq0}q^{2j}p
= \frac{p}{1 - q^2} = \frac{1}{1 + q} .
\]
Voor een eerlijke munt: $\frac1{1 + 1/2} = \frac23$.
(Verstandscontrole: oneven rangen moeten waarschijnlijker zijn,
want rang $1$ komt als eerste --- en inderdaad is
$\frac1{1+q} > \frac12$ altijd.)
\end{solution}

\begin{solution}{exo:b2:proba:10}
De \omterm{def:b2:proba:space}{gebeurtenissen} $B_N = \bigcap_{n=1}^N A_n^c$ dalen naar
$\bigcap_nA_n^c$, en wegens de \omterm{def:b2:proba:independence}{onafhankelijkheid} van de
complementen is $\P(B_N) = \prod_{n=1}^N(1 - p_n)$; de monotone
\omterm{def:b2:metric:continuity}{continuïteit} (\cref{thm:b2:proba:continuity}) geeft de getoonde
limiet. Logaritmen nemen: $\prod(1 - p_n) > 0$ dan en slechts dan
als $\sum-\ln(1 - p_n) < \infty$. Is $\sum p_n < \infty$, dan is
$p_n \to 0$ en $-\ln(1 - p_n) \sim p_n$: de logaritmereeks
convergeert. Is $\sum p_n = \infty$, dan dwingt $-\ln(1 - p_n) \geq
p_n$ divergentie af, dus is het product $0$. Dit stemt overeen met
Borel--Cantelli~2: voor $\sum p_n = \infty$ is niet alleen
$\P(\text{geen enkele }A_n\text{ treedt op}) = 0$, maar treden er
bijna zeker oneindig veel $A_n$ op.
\end{solution}

\begin{solution}{exo:b2:proba:11}
Zeg dat doosje $A$ het doosje is dat het eerst leeg wordt
aangetroffen, terwijl het andere er $k$ bevat. Dat betekent: onder
de eerste $2n - k$ grepen gingen er precies $n$ naar $A$ en $n - k$
naar $B$ (in een zekere volgorde), en greep nummer $2n - k + 1$
ging opnieuw naar $A$, dat leeg bleek. De grepen zijn
\omterm{def:b2:proba:independence}{onafhankelijke} eerlijke keuzes, dus heeft deze \omterm{def:b2:proba:space}{gebeurtenis} kans
$\binom{2n-k}{n}2^{-(2n-k)}\cdot\frac12$; verdubbelen (het lege
doosje kan er een van beide zijn) geeft
\[
\P(\text{het andere doosje bevat }k) =
\binom{2n-k}{n}\,2^{-(2n-k)} .
\]
Voor $n = 1$: $k = 1$ geeft $\binom11 2^{-1} = \frac12$ en $k = 0$
geeft $\binom21 2^{-2} = \frac12$: samen $1$, zoals het moet.
\end{solution}

\begin{solution}{exo:b2:proba:12}
(a) Is $\P(\{n\}) = c$ voor alle $n$, dan dwingt de
$\sigma$-additiviteit $1 = \sum_nc$ af: onmogelijk, of $c = 0$ is
(som $0$) of $c > 0$ (som oneindig). Er bestaat geen uniforme kans
op $\N$.

(b) Is $A \cap B = \emptyset$ en bestaan $d(A)$ en $d(B)$, dan is
$\abs{(A \sqcup B)\cap\intint1n} = \abs{A\cap\intint1n} +
\abs{B\cap\intint1n}$, dus $d(A \sqcup B) = d(A) + d(B)$: eindige
additiviteit op zulke paren. Elk singleton heeft een telfunctie die
uiteindelijk constant is, en dus dichtheid $0$, terwijl $d(\N^*) =
1$. Was $d$ $\sigma$-additief, dan zou $\N^* = \bigsqcup_k\{k\}$
geven dat $1 = \sum_k 0 = 0$: de dichtheid is eindig additief maar
niet $\sigma$-additief --- het axioma heeft inhoud.

(c) Zij $A = \bigcup_{k\geq0}\intint{4^k}{2\cdot4^k - 1}$ (blokken
van $4^k$ tot $2\cdot4^k - 1$). In $n = 2\cdot4^K - 1$ is de
telling $\sum_{k\leq K}4^k \sim \frac43 4^K$, wat een verhouding
$\to \frac23$ geeft; in $n = 4^{K+1} - 1$ is de telling
onveranderd, wat een verhouding $\to \frac13$ geeft. De verhouding
oscilleert tussen de limieten $\frac13$ en $\frac23$: geen
dichtheid.
\end{solution}

\begin{solution}{pb:b2:proba:1}
\textbf{1.} Een pad van lengte $n$ wordt bepaald door de
verzameling van zijn opwaartse stappen; in $k$ eindigen betekent
$u$ opwaartse en $n - u$ neerwaartse stappen met $u - (n - u) = k$,
dat wil zeggen $u = \frac{n+k}2$: mogelijk dan en slechts dan als
$n + k$ even is en $\abs k \leq n$, en wel op
$\binom{n}{(n+k)/2}$ manieren. Elk specifiek pad is één punt van de
eerlijke productmaat op $n$ worpen: kans $2^{-n}$. Bijgevolg is
$\P(S_n = k) = N_n(k)2^{-n}$.

\textbf{2.} $S_n$ heeft de pariteit van $n$, dus $S_{2n+1} \neq 0$;
en $u_n = N_{2n}(0)4^{-n} = \binom{2n}n4^{-n}$. Waarden: $u_1 =
\frac12$, $u_2 = \frac6{16} = \frac38$, $u_3 = \frac{20}{64} =
\frac5{16}$.

\textbf{3.} $\dfrac{u_n}{u_{n-1}} =
\dfrac{\binom{2n}n}{4\binom{2n-2}{n-1}} =
\dfrac{(2n)(2n-1)}{4n^2} = \dfrac{2n-1}{2n} < 1$: dalend. Volgens
\cref{ex:b2:comparison:centralbinomial} is $\binom{2n}n \sim
\frac{4^n}{\sqrt{\pi n}}$, dus $u_n \sim \frac1{\sqrt{\pi n}} \to
0$, en $\sum u_n$ divergeert door vergelijking met $\sum
n^{-1/2}$.

\textbf{4.} Gegeven een pad van $1$ naar $k$ dat $0$ raakt,
spiegel je zijn beginstuk (tot aan het \emph{eerste} bezoek aan
$0$) in de horizontale as: het resultaat is een pad van $-1$ naar
$k$, en de bewerking is een involutie --- elk pad van $-1$ naar $k
\geq 1$ moet $0$ kruisen, en zijn beginstuk terugspiegelen geeft
het oorspronkelijke terug. Bijgevolg zijn er $N_{n-1}(k + 1)$ paden
die raken (van $-1$ naar $k$ is de verplaatsing $k + 1$). Een pad
van $0$ naar $k$ dat na tijdstip $0$ boven $0$ blijft, begint met
een opwaartse stap en gaat dan in $n - 1$ stappen van $1$ naar $k$
zonder $0$ te raken: daarvan zijn er $N_{n-1}(k-1) -
N_{n-1}(k+1)$.

\textbf{5.} Met $m = \frac{n+k}2$ en $\binom{n-1}{m-1} = \frac
mn\binom nm$, $\binom{n-1}{m} = \frac{n-m}n\binom nm$:
\[
\frac{N_{n-1}(k-1) - N_{n-1}(k+1)}{N_n(k)}
= \frac{\binom{n-1}{m-1} - \binom{n-1}{m}}{\binom nm}
= \frac{m - (n - m)}{n} = \frac kn .
\]
Voor $n = 3$, $k = 1$: $N_3(1) = 3$ paden ($++-$, $+-+$, $-++$),
waarvan alleen $++-$ positief blijft ($+-+$ keert op tijdstip $2$
naar $0$ terug): één op drie, en $\frac kn = \frac13$.

\textbf{6.} Wegens de symmetrie is de kans $2\P(S_i > 0\ \forall i
\leq 2n)$. Sommeren over het eindpunt $2k$ en vraag 4 gebruiken
(met $n$ vervangen door $2n$) geeft
\[
\P(S_i > 0\ \forall i) = 2^{-2n}\sum_{k\geq1}
\bigl(N_{2n-1}(2k-1) - N_{2n-1}(2k+1)\bigr)
= 2^{-2n}\,N_{2n-1}(1),
\]
een telescoperende som. Nu is $N_{2n-1}(1) = \binom{2n-1}{n}$ en
$2\binom{2n-1}n = \binom{2n}n$ (Pascal), dus is de getoonde kans
$2\cdot2^{-2n}\binom{2n-1}n = \binom{2n}n4^{-n} = u_n$.

\textbf{7.} De \omterm{def:b2:proba:space}{gebeurtenissen} $D_n = \{S_i \neq 0,\ i \leq 2n\}$
dalen, met als doorsnede ``nooit een terugkeer''; volgens de
monotone \omterm{def:b2:metric:continuity}{continuïteit} en vraag 6 is $\P(\text{geen terugkeer}) =
\lim u_n = 0$: de wandeling keert bijna zeker terug. Bovendien is
$f_n = \P(D_{n-1}) - \P(D_n) = u_{n-1} - u_n$, en volgens vraag 3
\[
u_{n-1} - u_n = u_n\Bigl(\frac{2n}{2n-1} - 1\Bigr) =
\frac{u_n}{2n-1};
\qquad
\sum_{n\geq1}f_n = u_0 - \lim u_n = 1 .
\]

\textbf{8.} $2n\,f_n = \frac{2n}{2n-1}u_n \geq u_n$, en $\sum u_n =
\infty$ (vraag 3): de reeks $\sum 2nf_n$ divergeert. De eerste
terugkeer is zeker maar heeft geen eindige gemiddelde wachttijd ---
de wandeling is \emph{nulrecurrent}, in de woordenschat die
\cref{ch:b2:randomvar} zal leveren.

\textbf{9.} De \omterm{def:b2:proba:space}{gebeurtenis} ``minstens $k$ terugkeren'' is de
disjuncte \omterm{def:b2:structures:countable}{aftelbare} vereniging, over $0 < n_1 < \dots < n_k$, van
de \omterm{def:b2:proba:space}{gebeurtenissen} ``de eerste $k$ terugkeren gebeuren precies op de
tijdstippen $2n_1, \dots, 2n_k$''. Zulk een \omterm{def:b2:proba:space}{gebeurtenis} is de
doorsnede van $k$ \omterm{def:b2:proba:space}{gebeurtenissen} die van de disjuncte blokken
worpen $\intint1{2n_1}$, $\intint{2n_1+1}{2n_2}$, \dots afhangen,
waarbij elk blok van een verse wandeling vraagt dat zij haar eerste
terugkeer precies na het toegemeten aantal stappen maakt; wegens de
\omterm{def:b2:proba:independence}{onafhankelijkheid} van de blokken is haar kans
$f_{n_1}f_{n_2-n_1}\cdots f_{n_k-n_{k-1}}$. Sommeren in pakketten
(\cref{ch:b2:series}, alle termen niet-negatief) geeft
\[
\P(\text{minstens }k\text{ terugkeren})
= \Bigl(\sum_{n\geq1}f_n\Bigr)^{\!k} = 1^k = 1 .
\]
De \omterm{def:b2:proba:space}{gebeurtenissen} dalen in $k$, dus geeft de monotone \omterm{def:b2:metric:continuity}{continuïteit}
dat $\P(\text{oneindig veel terugkeren}) = 1$: recurrentie.

\textbf{10.} Volgens vraag 9 maakt de wandeling oneindig veel
uitstapjes weg van $0$. De eerste stap van elk uitstapje is een
verse munt, \omterm{def:b2:proba:independence}{onafhankelijk} van alles ervoor: de kans dat de eerste
$m$ uitstapjes alle neerwaarts beginnen, is $2^{-m}$. Om $1$ te
bereiken heeft de wandeling maar één opwaarts beginnend uitstapje
nodig (vanuit $<0$ moet zij door $0$ voordat zij $1$ bereikt, want
de stappen zijn $\pm1$), dus is $\P(\text{nooit }1\text{ raken})
\leq 2^{-m}$ voor elke $m$: de wandeling raakt $1$ bijna zeker.
Ontbinden we over de (bijna zeker eindige) raaktijd, dan is de daar
herstarte wandeling een verse wandeling die in $1$ begint: met
inductie raakt zij bijna zeker elke $k \geq 1$, en wegens de
symmetrie elke $k \leq -1$. Herstarten wij ten slotte bij het
eerste bezoek aan $k$, dan is vraag 9 van toepassing op de verse
wandeling: elke plaats wordt bijna zeker oneindig vaak bezocht.

\textbf{11.} De \omterm{def:b2:proba:space}{gebeurtenissen} $A_n = \{S_{2n} = 0\}$ zijn verre
van \omterm{def:b2:proba:independence}{onafhankelijk} (in $0$ zijn op tijdstip $2n$ maakt in $0$ zijn
op tijdstip $2n + 2$ veel waarschijnlijker dan $u_{n+1}$), dus is
Borel--Cantelli~2 niet beschikbaar, en het hele werk van Deel II
was dan ook haar te vervangen. De andere richting heeft geen
\omterm{def:b2:proba:independence}{onafhankelijkheid} nodig: \emph{als} $\sum\P(A_n)$ convergeert, dan
levert Borel--Cantelli~1 bijna zeker eindig veel terugkeren. Die
implicatie is de motor van elk bewijs van transiëntie hieronder.

\textbf{12.} Een terugkeer op tijdstip $2n$ vergt $n$ opwaartse en
$n$ neerwaartse stappen: $\P(S_{2n} = 0) = \binom{2n}np^nq^n =
u_n(4pq)^n$, en $4pq = 1 - (p - q)^2 < 1$ voor $p \neq \frac12$.
Omdat $u_n \leq 1$, wordt de reeks $\sum\P(S_{2n} = 0)$ gedomineerd
door de meetkundige $\sum(4pq)^n$: convergent. Volgens
Borel--Cantelli~1 is $\P(S_{2n} = 0 \text{ oneindig vaak}) = 0$:
bijna zeker eindig veel terugkeren.

\textbf{13.} Voor $n + k$ even is $\P(S_n = k) =
\binom{n}{\frac{n+k}2}p^{\frac{n+k}2}q^{\frac{n-k}2}$; de
binomiaalcoëfficiënt is hoogstens de centrale, en
$p^{\frac{n+k}2}q^{\frac{n-k}2} = (pq)^{n/2}(p/q)^{k/2}$, wat de
gestelde grens $\leq 2^n(pq)^{n/2}(p/q)^{k/2} =
(4pq)^{n/2}(p/q)^{k/2}$ geeft, \omterm{def:b2:series:summable}{sommeerbaar} in $n$ omdat
$\sqrt{4pq} < 1$. Borel--Cantelli~1: plaats $k$ wordt bijna zeker
eindig vaak bezocht; de vereniging over $k \in \Z$ van de
uitzonderlijke nulgebeurtenissen is nog altijd nul (\omterm{def:b2:structures:countable}{aftelbare}
subadditiviteit). Bijna zeker wordt elke plaats eindig vaak
bezocht, dus verlaat de rij gehele getallen $(S_n)$ elk begrensd
venster voorgoed: $\abs{S_n} \to \infty$.

\textbf{14.} $\P(S_n \neq 0,\ 1 \leq n \leq 200) = u_{100} =
\binom{200}{100}4^{-100} \approx \frac1{\sqrt{100\pi}} \approx
0.056$: meer dan één kans op twintig dat $200$ eerlijke worpen
nooit een gelijkstand geven. Het verval $1/\sqrt{\pi n}$ is
tergend traag: de zekerheid van een gelijkstand (vraag 7) is
verenigbaar met heel lange stukken zonder gelijkstand --- een
eerste voorproefje van de arcsinusverschijnselen van Deel IV.

\textbf{15.} Stel voorwaardelijk op de eerste stap. Is $X_1 = +1$,
dan is $T_1 = 1$, en $f_1 = \frac12$ klopt. Is $X_1 = -1$, dan moet
de wandeling van $-1$ naar $1$ klimmen; volgens de ontbinding in
blokken splitst een eerste terugkeer naar $0$ op tijdstip $2n$ als:
één stap omlaag, en dan een verse wandeling die vanuit $-1$ voor
het eerst $0$ bereikt --- gelijkwaardig: een verse wandeling die
voor het eerst $+1$ bereikt --- in $2n - 1$ stappen, of de
symmetrische \omterm{def:b2:proba:space}{gebeurtenis} omhoog. Beide tekens dragen evenveel bij:
\[
f_n = 2\cdot\tfrac12\,\P(T_1 = 2n - 1) = \P(T_1 = 2n-1) .
\]
Bijgevolg is $\P(T_1 < \infty) = \sum f_n = 1$, terwijl $\sum_n(2n
- 1)f_n = \sum_n u_n = \infty$ volgens vraag 7: de wandeling
bereikt $1$ bijna zeker, in oneindige gemiddelde tijd.

\textbf{16.} Partitioneer $\{M_n \geq k\}$ naar de eindwaarde $S_n
= m$. Voor $m \geq k$ is de voorwaarde $M_n \geq k$ automatisch.
Voor $m < k$ spiegel je het pad na zijn \emph{eerste} bezoek aan
niveau $k$: dit is een bijectie tussen $\{M_n \geq k, S_n = m\}$ en
$\{S_n = 2k - m\}$ (elk pad dat in $2k - m > k$ eindigt, bezoekt
$k$; terugspiegelen is de inverse). Bijgevolg is
\[
\P(M_n \geq k)
= \sum_{m > k}\P(S_n = m) + \P(S_n = k)
+ \sum_{m < k}\P(S_n = 2k - m)
= 2\P(S_n > k) + \P(S_n = k).
\]

\textbf{17.} Op het even tijdstip $2n$ met $k = 1$: $\P(S_{2n} =
1) = 0$ en $\P(S_{2n} > 1) = \P(S_{2n} \geq 2)$, dus
\[
\P(M_{2n} \geq 1) = 2\P(S_{2n} \geq 2)
= \P(S_{2n} \geq 2) + \P(S_{2n} \leq -2)
= 1 - u_n .
\]
Dus $\P(S_i \leq 0\ \forall i \leq 2n) = u_n$: de wandeling staat
in de eerste $2n$ stappen precies even vaak nooit voor als zij
nooit gelijk staat (vraag 6) --- twee heel verschillende
\omterm{def:b2:proba:space}{gebeurtenissen}, gedragen door dezelfde $u_n$.

\textbf{18.} $\{L_{2n} = 2k\} = \{S_{2k} = 0\} \cap \{\text{de
wandeling van de worpen } 2k+1, \dots, 2n \text{ heeft geen
nulpunt}\}$. De twee \omterm{def:b2:proba:space}{gebeurtenissen} hangen van disjuncte blokken
worpen af en zijn dus \omterm{def:b2:proba:independence}{onafhankelijk}; de eerste heeft kans $u_k$, de
tweede $u_{n-k}$ volgens vraag 6 toegepast op de verse wandeling
van $2n-2k$ stappen. Bijgevolg is $\P(L_{2n} = 2k) = u_ku_{n-k}$.
Omdat $L_{2n}$ precies de waarden $0, 2, \dots, 2n$ aanneemt,
sommeren deze kansen tot $1$: $\sum_{k=0}^nu_ku_{n-k} = 1$, een
binomiale identiteit geleverd door een kansrekenkundige partitie.

\textbf{19.} De symmetrie is onmiddellijk: $u_ku_{n-k} =
u_{n-k}u_k$. Omdat $u_j$ daalt in $j$, is het product
$u_ku_{n-k}$ het kleinst voor centrale $k$ en het grootst in de
uitersten $k \in \{0, n\}$, waar het gelijk is aan $u_n$;
kwantitatief is $u_ku_{n-k} \approx \frac1{\pi\sqrt{k(n-k)}}$ in
het midden, tegenover $u_n \approx \frac1{\sqrt{\pi n}}$ aan de
randen. Voor $n = 5$: $\P(L_{10} = 0) = \P(L_{10} = 10) = u_5 =
\frac{63}{256} \approx 0.246$, terwijl $\P(L_{10} = 4) = u_2u_3 =
\frac38\cdot\frac5{16} = \frac{15}{128} \approx 0.117$. In een lang
eerlijk spel valt de laatste gelijkstand het waarschijnlijkst
vlakbij het begin of vlakbij het einde: één speler staat meestal
gedurende enorme stukken voor, zonder enige scheefheid in de munt.

\textbf{20.} Het beeld: op tijdstip $n$ leeft de wandeling op de
schaal $\sqrt n$ (de binomiale spreiding van vraag 3 --- $u_n \sim
1/\sqrt{\pi n}$ is de hoogte van de centrale piek); zij keert met
kans $1$ oneindig vaak naar $0$ terug (Deel II), en toch heeft de
wachttijd tussen de terugkeren een divergent gemiddelde (vraag 8),
en daarom kunnen afzonderlijke uitstapjes een positief aandeel van
elke horizon innemen; navenant is de laatste gelijkstand van een
spel met $2n$ stappen uitgesmeerd met de uiterste waarden als de
waarschijnlijkste (vragen 18--19), en heeft nooit voorstaan
dezelfde traag vervallende kans $u_n$ als nooit gelijk staan (vraag
17). Zekerheid in de limiet, volharding op elke eindige horizon:
dat is de eerlijke wandeling.

\textbf{21.} Partitioneer $\{S_{2n} = 0\}$ ($n \geq 1$) naar de
eerste terugkeertijd $2k$, $1 \leq k \leq n$: het eerste blok van
$2k$ worpen verwezenlijkt een eerste terugkeer, de overige $2n -
2k$ worpen verwezenlijken een terugkeer van een verse wandeling, en
de blokken zijn \omterm{def:b2:proba:independence}{onafhankelijk}: $u_n = \sum_{k=1}^nf_ku_{n-k}$.
Beide reeksen $U(x) = \sum u_nx^n$ en $F(x) = \sum f_nx^n$ hebben
straal $\geq 1$ (coëfficiënten in $\intcc01$), en het
\omterm{thm:b2:series:fubini}{Cauchy-product} (\cref{ch:b2:powerseries}) geeft voor $0 \leq x < 1$
\[
U(x) - 1 = \sum_{n\geq1}\Bigl(\sum_{k=1}^n
f_ku_{n-k}\Bigr)x^n = F(x)\,U(x),
\qquad\text{dat wil zeggen}\qquad
U(x)\bigl(1 - F(x)\bigr) = 1 .
\]

\textbf{22.} Als $x \uparrow 1$ stijgen $U(x)$ en $F(x)$
(niet-negatieve coëfficiënten); elke partiaalsom $\sum_{n\leq
N}u_n$ is een limiet van $\sum_{n\leq N}u_nx^n \leq U(x)$, dus
$U(x) \uparrow \sum u_n \in \intoc0{+\infty}$, en evenzo $F(x)
\uparrow f = \sum f_n$. Is $\sum u_n = \infty$, dan $1 - F(x) =
1/U(x) \to 0$, dus $f = 1$. Is $\sum u_n = S < \infty$, dan $1 - f
= 1/S > 0$, dus $f < 1$. Controles: eerlijke wandeling, $\sum u_n =
\infty$ en $f = 1$ (vragen 3, 7); scheve wandeling, $\sum
u_n(4pq)^n < \infty$ en navenant $f = 1 -
1/\sum_{n\geq0}u_n(4pq)^n < 1$, in overeenstemming met de
bijna-zekere eindigheid van het aantal terugkeren (vraag 12).

\textbf{23.} Voor de vier stappen $(\pm1, 0), (0, \pm1)$ van de
wandeling op $\Z^2$ zijn de aangroeiingen van $U = X + Y$ en $V = X
- Y$: $(+,+)$ voor $(1,0)$, $(+,-)$ voor $(0,1)$, $(-,+)$ voor
$(0,-1)$ en $(-,-)$ voor $(-1,0)$ --- elk tekenpaar met kans
$\frac14 = \frac12\cdot\frac12$: de twee coördinaatwandelingen
$(U_n)$ en $(V_n)$ zijn \omterm{def:b2:proba:independence}{onafhankelijke} eerlijke wandelingen op
$\Z$. Omdat $S^{(2)}_{2n} = (0,0)$ dan en slechts dan als $U_{2n} =
0$ en $V_{2n} = 0$, is
\[
\P\bigl(S^{(2)}_{2n} = (0,0)\bigr) = u_n^2 \sim \frac1{\pi
n}, \qquad \sum_nu_n^2 = \infty .
\]
De vernieuwingsidentiteit van vraag 21 en de tweedeling van vraag
22 gebruikten niets eendimensionaals (alleen de ontbinding over de
eerste terugkeer en de \omterm{def:b2:proba:independence}{onafhankelijkheid} van disjuncte blokken),
dus geeft $\sum u_n^{(2)} = \infty$ dat $f^{(2)} = 1$, en het
argument van vraag 9 verscherpt het: de wandeling op $\Z^2$ keert
bijna zeker oneindig vaak naar de oorsprong terug.

\textbf{24.} Met de aangenomen grens $\P(S^{(3)}_{2n} = 0) \leq
Cn^{-3/2}$ convergeert de reeks, en Borel--Cantelli~1 geeft bijna
zeker eindig veel terugkeren: de wandeling op $\Z^3$ is transiënt
(en dezelfde grens met exponent $-d/2$ handelt elke $d \geq 3$ af).
Alles samen: de \emph{stelling van Pólya} --- de \omterm{pb:b2:proba:1}{eenvoudige
toevalswandeling} is recurrent op $\Z$ en $\Z^2$, en transiënt op
$\Z^d$ voor $d \geq 3$. Een dronken man vindt de weg naar huis; een
dronken vogel misschien niet.

\textbf{25.} Het tellen van paden en de spiegeling brachten de
exacte verdelingen voort ($u_n$, de stemmenstelling, $f_n$, het
maximum, het laatste nulpunt); de monotone \omterm{def:b2:metric:continuity}{continuïteit} zette elke
uitspraak over een limiet (``keert minstens eenmaal terug'',
``oneindig vaak'') om in een limiet van kansen met eindige horizon;
de \omterm{def:b2:proba:independence}{onafhankelijkheid} van disjuncte blokken dreef de
vernieuwingsontbindingen aan (vragen 9, 18, 21) --- zij is het
\omterm{def:b2:structures:countable}{aftelbare} skelet van de markov-eigenschap; Borel--Cantelli~1 was
het wapen voor de transiëntie (vragen 12--13, 24), zonder
\omterm{def:b2:proba:independence}{onafhankelijkheid} nodig te hebben; en de vernieuwingsidentiteit
ordende alles tot de tweedeling $\sum u_n = \infty \iff$
recurrentie. Het enige analytische ingrediënt is de lokale
schatting $u_n \sim 1/\sqrt{\pi n}$: haar kwadraat $1/(\pi n)$
divergeert nog altijd (dimensie $2$, recurrent), terwijl
$n^{-3/2}$ convergeert (dimensie $3$, transiënt) --- de stelling
van Pólya is uiteindelijk een uitspraak over de divergentie van
$\sum n^{-d/2}$.
\end{solution}
